\documentclass[../report.tex]{subfiles}

\begin{document}

\subsection{Linux - Complete Pentest (Mesa)}

\subsubsection{OSINT}

All publicly available information relevant to security concerns will be listed in this section.
The whole writeup created during the penetration test can be viewed in the appended files under:\\
Assignment 2 - Complete Pentest (Mesa) $\longrightarrow$ 1. OSINT $\longrightarrow$ osint.txt

\textbf{Marquis Buckerzerg (CEO)}\\[0.2em]
\begin{tabular}{p{9.3cm} | l}
- Full name: \textit{Marquis Buckerzerg} & [related to their usernames]\\ \hline
- Loves their dog \enquote{tinkerbell}, calls them their \enquote{sunshine} & [related to their passwords]
\end{tabular}
\begin{figure}[!htb]
	\centering
	
\includegraphics[width=\linewidth]{pentest/1_active_scan/1_osint/dog}
	\caption{X.com post by M. Buckerzerg about his dog tinkerbell}
	\vspace{-2mm}
	\caption*{\small Source: \url{https://x.com/MBuckerzerg/status/1494491183372017668}}
\end{figure}

\textbf{Samantha Groves (Database technician)}\\[0.2em]
\begin{tabular}{p{9.3cm} | l}
- Last name: \textit{Groves} and nickname: \textit{sam} & [related to their usernames]\\ \hline
- Uses the word \enquote{machine} in their team description & [related to their passwords]
\end{tabular}

\begin{figure}[!htb]
	\centering
	
\includegraphics[width=0.5\linewidth]{pentest/1_active_scan/1_osint/sam}
	\caption{Mesa homepage - Team description of Samantha Groves}
\end{figure}

\textbf{Gordon Freeman (Business partner)}\\
\begin{tabular}{p{9.3cm} | l}
- Last name: \textit{Freeman} & [related to their usernames]\\ \hline
\end{tabular}

\begin{figure}[!htb]
	\centering
	
\includegraphics[width=0.3\linewidth]{pentest/1_active_scan/1_osint/freeman}
	\caption{Mesa homepage - Listing of Gordon Freeman}
\end{figure}

\newpage

\subsubsection{Active Scan / Active Interaction}
\paragraph{Banner Grabbing / Fingerprinting}
\subparagraph{Nmap}
Nmap was used to scan the target system with the following parameters:

\begin{figure}[!htb]
	\centering
	\includegraphics[width=\linewidth]{pentest/1_active_scan/2_nmap/params}
	\caption{Nmap parameters}
\end{figure}

The reasoning for each chosen parameter is the same as it was during the Buffer Overflow assignment (see \fullref{bo:scan-setup}). The scripts were omitted since other methods of vulnerability scanning were employed later.
The nmap scan yielded the following results:

\begin{tabular}{l | l | l}
\textbf{Open Ports} & \textbf{Running Service} & \textbf{Technology}\\
22 & Secure Shell & OpenSSH 8.4p1 Debian 5 (protocol 2.0)\\ \hline
80 & Webserver 1 & nginx 1.18.0\\ \hline
1337 & Webserver 2 & nginx 1.18.0\\ \hline 
5355 & Link-Local Multicast Name Resolution & LLMNR\\ \hline 
8080 & Webserver 3 & nginx 1.18.0
\end{tabular}

Nmap discovered that the machine was accessible via SSH on the standard port 22, and that the machine does local name resolution via the LLMNR protocol on port 5355.
%TODO check https://www.syss.de/pentest-blog/llmnr-das-oft-vergessene-einfallstor-ins-netzwerk
LLMNR is a protocol, which allows machines in the local network to resolve requests locally if DNS were to fail \autocite{rfc4795}. 
There are three webservers running on the machine, on port 80, 1337, and 8080 respectively, reachable via HTTP.

\subparagraph{Netcat}
Connection attempts were made via netcat.\\
The SSH port allowed connecting, but since no credentials or keys were known at this point, further access was not possible (see \cref{nc22}).
\begin{figure}[!htb]
	\centering
	\includegraphics[width=\linewidth]{pentest/1_active_scan/3_banner_grabbing/nc/nc22.png}
	\caption{Banner with nc grabbing on port 22}
	\label{nc22}
\end{figure}
Other banner grabbing attempts yielded no further information, since no service returned a substantive response (see \cref{nc_others}). 
\begin{figure}[!htb]
	\centering
	\includegraphics[width=0.7\linewidth]{pentest/1_active_scan/3_banner_grabbing/nc/nc_other}
	\caption{Banner grabbing with nc on ports 80, 1337, 5355, 8080}
	\label{nc_others}
\end{figure}

\newpage

\subparagraph{Curl}
Information about the webservers was gathered via curl.
The servers on port 80 and 1337 seemed to be configured similarly.
Both allow connection via HTTP, returning http/text type content, and don't seem to have any additional headers set up (see \cref{curl80,curl1337}).

\begin{figure}[!htb]
	\centering
	\includegraphics[width=0.8\linewidth]{pentest/1_active_scan/3_banner_grabbing/curl/80}
	\caption{Fingerprinting with curl on port 80}
	\label{curl80}
\end{figure}

\begin{figure}[!htb]
	\centering
	\includegraphics[width=0.8\linewidth]{pentest/1_active_scan/3_banner_grabbing/curl/1337}
	\caption{Fingerprinting with curl on port 1337}
	\label{curl1337}
\end{figure}

They were last modified only seconds apart, indicating they might contain similar vulnerabilities, if they were to be discovered.
Since no additional header seems to be configured, the webpages might be vulnerable to clickjacking, cross-site-scripting (XSS), and a host of other attacks \autocite{owasp:headers}.


\newpage

\paragraph{Directory Fuzzing}

\subparagraph{Ffuf}
\begin{figure}[!htb]
	\centering
	\includegraphics[width=\linewidth]{pentest/1_active_scan/4_fuzzing/ffuf/params}
	\caption{Ffuf parameters}
	\label{dirbuster_complete_setup}
\end{figure}

The parameters for Ffuf were chose with the following reasoning \autocite{ffuf:wiki}:
\begin{itemize}
\renewcommand\labelitemi{-}
\itemBF{c}
Colorizes the ouput for easier readability.
\itemBF{ac}
Automatically calibrates filtering options for better efficiency.
\item{\textbf{e} \textit{.php,.html,.txt,.md,.sql,.save,.backup,.js,.jsx}}\\
Sets the file extensions Ffuf fuzzes for. Passed a list of possibly interesting file formats.
\item{\textbf{w} \textit{/usr/share/seclists/Discovery/Web-Content/common.txt}}\\
Uses a wordlists included in the Kali Linux distribution, containing common directory and file names.
\item{\textbf{-u} \textit{http://mesa/FUZZ}}\\
Sets the target url to scan the Mesa homepage. \textit{/FUZZ} indicates that Ffuf should fuzz from the base directory onwards.
\item[--]{\textbf{recursion-depth} \textit{5}}\\
Sets the recursion depth to 5. Higher recursion depths didn't seem to yield any more interesting results.
\itemBF{v}
Sets the output to verbose mode to enable following the progress on standard output.
\item{\textbf{o} \textit{ffuf\_80.json}}\\
Saves the output to the file ffuf\_80.json. The suffix 80 was chosen to distinguish it from the output for the scan on port 1337.
\end{itemize}

The results for port 80 will be discussed in the following section, since dirbuster output is more readable.
The complete output of the Ffuf scan can be viewed in the appended files under:\\
{\small
\textit{Assignment 2 - Complete Pentest (Mesa) $\longrightarrow$ 2. Active Scan / Active Interaction $\longrightarrow$ ffuf $\longrightarrow$ ffuf\_80.png'}\\
and \\
{\small
\textit{Assignment 2 - Complete Pentest (Mesa) $\longrightarrow$ 2. Active Scan / Active Interaction $\longrightarrow$ ffuf $\longrightarrow$ ffuf\_80.json'}}\\
}
or in the appendix of this document (see \cref{ffuf_80_full}).

The scan for port 1337 did not reveal any information of note. The base directory only seems to contain the \textit{index.html} file, with no other files of interest (see \cref{ffuf_1337}).

\begin{figure}[!htb]
	\makebox[1.03\linewidth][c]{%
		\includegraphics[width=1.2\linewidth]{pentest/1_active_scan/4_fuzzing/ffuf/ffuf_1337}
	}
	\caption{Ffuf fuzzing on port 1337}
	\label{ffuf_1337}
\end{figure}
The parameters used were exactly the same as before when scanning port 80, except that the parameter \textbf{- u} now received the url http://mesa:1337/FUZZ, to scan on the respective port.

It can be assumed that the webpage is a static html file without much interactability. This will be confirmed during the examination of the webpage later during browsing.
This webpage has been examined with dirbuster as well, which did not yield any different results. The results can be viewed in the appendix to this document for completeness sake. (see \cref{dirbuster_1337_setup,dirbuster_1337_result}).

\subparagraph{Dirbuster}
To scan port 80 with dirbuster, the following parameters were chosen (see \cref{dirbuster_80_setup}):

\begin{tabular}{l | l}
\textbf{Target URL} & \textbf{http://mesa}\\
\tabnote{The target was the target machine. Supplying a port was not necessary, it defaults to port 80.}
\textbf{Work Method} & \textbf{Auto Switch (HEAD and GET)}\\
\tabnote{Automatically switches between HEAD and GET for better scanning efficiency.}
\textbf{Number of threads} & \textbf{500}\\
\tabnote{Set the thread count to the maximum possible number to achieve the fastest possible speed.}
\textbf{Scanning type} & \textbf{List based brute force}\\
\tabnote{Select directories and files to check from the supplied wordlist (same as with Ffuf).}
\textbf{Starting options} & \textbf{Brute force directories and files, be recursive, start from /}\\
\tabnote{From the base directory, recursively search for directories and files.}
\end{tabular}

\begin{figure}[!htb]
	\centering
	\includegraphics[width=\linewidth]{pentest/1_active_scan/4_fuzzing/dirbuster/setup/80-setup}
	\caption{Dirbuster setup}
	\label{dirbuster_80_setup}
\end{figure}

\begin{figure}[!htb]
	\centering
	\includegraphics[width=\linewidth]{pentest/1_active_scan/4_fuzzing/dirbuster/80-result}
	\caption{Dirbuster result - Directory tree of 192.168.61.65:80}
	\label{dirbuster_80_setup}
\end{figure}
Dirbuster discovered the following directories / files of note:

\begin{itemize}
\itemBF{backup.sql}
A database backup was discovered inside an indexed backup directory \textit{/backup}. 
\itemBF{login.php \& management.php}
There are two php pages \textit{login} and \textit{management} not directly accessible via the homepage.
Judging by the file names they might provide an entry point to get to sensitive data.
\end{itemize}

\newpage
\paragraph{Technology discovery}
\subparagraph{Whatweb}
To scan the technologies used by the webservers, the file \textit{urls.txt} was created.
It contained the following content:
\vspace{-6mm}
\begin{verbatim}
http://192.168.61.65:80
http://192.168.61.65:1337
http://192.168.61.65:8080
\end{verbatim}

\begin{figure}[!htb]
	\centering
	\includegraphics[width=\linewidth]{pentest/1_active_scan/5_technologies/params.png}
	\caption{Whatweb parameters}
	\label{whatweb_params}
\end{figure}

The chosen parmeters (see \cref{whatweb_params}) were selected with the following reasoning \autocite{whatweb:man}:
\begin{itemize}
\renewcommand\labelitemi{\texttt{-{-}}}
\item{\textbf{log-verbose}=whatweb.txt}\\
Write the ouput to a file \textit{whatweb.txt} for later analysis.
\item{\textbf{input-file}=urls.txt}\\
Use the urls provided in the file \textit{urls.txt} to scan multiple targets at once.
\item[-]{\textbf{v}}
Sets output to verbose to be able to follow results in the terminal emulator.
\end{itemize}


\newpage
\begin{figure}[!htb]
	\centering
	\includegraphics[width=\linewidth]{pentest/1_active_scan/5_technologies/1337}
	\caption{Whatweb results for port 1337}
	\label{whatweb_1337}
\end{figure}
\newpage
\begin{figure}[!htb]
	\centering
	\makebox[\linewidth][c]{%
		\includegraphics[width=0.82\linewidth, height=21.4cm]{pentest/1_active_scan/5_technologies/80}
	}
	\caption{Whatweb results for port 80}
	\label{whatweb_80}
\end{figure}
\newpage

\paragraph{Vulnerability scan}
\subparagraph{Nikto}
\begin{figure}[!htb]
	\centering
	\includegraphics[width=\linewidth]{pentest/1_active_scan/6_vulns/params}
	\caption{Nikto parameters}
	\label{nikto_params}
\end{figure}

The selected parameters (see \cref{nikto_params}) were chosen with the following reasoning \autocite{nikto:params}:

\newpage
\subsubsection{Examining discovered web pages}
The page \textit{login.php} presented a login form asking for a username and a password (see \cref{mesa_login}).
\begin{figure}[!htb]
	\centering
	\includegraphics[width=0.8\linewidth]{pentest/2_login/login_page.png}
	\caption{Mesa login page}
	\label{mesa_login}
\end{figure}

Trying to access management.php directly redirected to the login page (see \cref{mesa_redirect}). A login was necessary to access the protected area.

\begin{figure}[!htb]
	\centering
	\includegraphics[width=0.9\linewidth]{pentest/2_login/redirect}
	\caption{Redirect to Mesa login page without logging in}
	\label{mesa_redirect}
\end{figure}

The \textit{/backup} directory was accessible (see \cref{backup_page}), and the previously discovered SQL backup was downloaded.
\begin{figure}[!htb]
	\centering
	\includegraphics[width=\linewidth]{pentest/2_login/backup_page}
	\caption{Acessing the backup directory in the browser}
	\label{backup_page}
\end{figure}

The backup contained web-app usernames, password hashes, information about the user's admin status, if they own a unix account, unix usernames, and salary information (see \cref{backup_data}).
\begin{figure}[!htb]
	\centering
	\includegraphics[width=\linewidth]{pentest/2_login/backup_data}
	\caption{Acessing the SQL backup}
	\label{backup_data}
\end{figure}

At this point, enough information has been gathered to make educated guesses about credentials.
Guessing the words Mr. Buckerzerg seems to ascribe a high emotional value to lead to a hit:
\enquote{sunshine} worked as a password for the management system account \enquote{marq} (see \cref{access_marq}).
\begin{figure}[!htb]
	\centering
	\includegraphics[width=\linewidth]{pentest/2_login/access_marq}
	\caption{Acessing the management system as user marq}
	\label{access_marq}
\end{figure}
Checking the website revealed that the information from the database backup is also displayed in the management system.
Reading the system log, we can see that the user freeman changed his password (see \cref{syslog}). 
\begin{figure}[!htb]
	\centering
	\includegraphics[width=\linewidth]{pentest/2_login/syslog}
	\caption{System log showing freeman changed their password}
	\label{syslog}
\end{figure}

The hash displayed in the frontend does not match the has from the backup (see \cref{freeman_pw_change}).:
\begin{figure}[!htb]
	\centering
	
\includegraphics[width=\linewidth]{pentest/2_login/freeman_pw_change}
	\caption{System log showing freeman changed their password}
	\label{freeman_pw_change}
\end{figure}

\begin{tabular}{l l}
Old Hash: & fddd21b9d7ce17da93c30fa5a653a1df\\
New Hash: & df64dc2eb4a0b85091dd31eb4923eaac
\end{tabular}

It is feasible to assume that the password \enquote{machine} for the user groves coul be guessed as well, based on their description on the Mesa homepage. Since the password hashes were known at this point, the credential were able to be extracted without guessing as well. The hashes will be cracked at a later point with Linux tools, but it is also possible to get the credentials via online hash cracking services (see \cref{frankfurt}).
\begin{figure}[!htb]
	\centering
	\includegraphics[width=\linewidth]{pentest/2_login/frankfurt}
	\caption{Cracking password hashes via crack station}
	\vspace{-2mm}
	\caption*{\small Source: \url{https://crackstation.net}}
	\label{frankfurt}
\end{figure}

The publicly available information therefore lead to the compromise of all credentials for the web application:

\begin{tabular}{p{4cm} | p{5cm}}
\textbf{Username} & \textbf{Password}\\ \hline
marq & sunshine\\ \hline
freeman & ??????\\ \hline
groves & machine\\ \hline
joe & hackerman\\ \hline
heuzeroth & anwendungssicherheit
\end{tabular}
\newpage

\subsubsection{Logging into the Web-Application}
\end{document}